% GENERATED by scripts/make_prompt_appendix.py -- do not edit by hand.
% Regenerate after any change to diagnostics/prompts_L*.json.

\section{Prompt formulations}
\label{app:prompts}

The caption ablation (\S\ref{sec:results:ablation}) varies only the text embedding. All five prompt sets are reproduced below verbatim, generated directly from the files the ablation consumed.

\subsection*{\texttt{bladder}}
\begin{description}[leftmargin=3.2em,style=nextline]
\item[\textbf{L1}] bladder
\item[\textbf{L2}] Axial CT scan of the bladder.
\item[\textbf{L3}] The urinary bladder is located in the top half of the image at the horizontal center, positioned symmetrically across the vertical midline.
\item[\textbf{L4}] Axial CT scan of the pelvis showing the urinary bladder as a smooth-walled fluid-filled hypodense oval. It sits at the horizontal center of the image in the top half, symmetric across the vertical midline. The rectum appears as a smaller structure at the bottom of the image behind it. The pubic symphysis is visible at the very top image edge in front of it. Wall thickness is thin and uniform.
\item[\textbf{L5}] Axial CT scan of the pelvis showing the urinary bladder as a smooth-walled fluid-filled hypodense oval. It sits at the horizontal center of the image in the bottom half, symmetric across the vertical midline. The rectum appears as a smaller structure at the top of the image behind it. The pubic symphysis is visible at the very bottom image edge in front of it. Wall thickness is thin and uniform.
\end{description}

\subsection*{\texttt{femur-left}}
\begin{description}[leftmargin=3.2em,style=nextline]
\item[\textbf{L1}] femur-left
\item[\textbf{L2}] Axial CT scan of the femur left.
\item[\textbf{L3}] The left femur is located on the right side of the image.
\item[\textbf{L4}] Axial CT scan of the thigh showing the left femur in transverse cross-section. Appears on the RIGHT side of the image as a bright hyperdense cortical ring surrounding lower-density cancellous bone with a central medullary canal. Adductor muscles border it toward the image center (medially). Vastus lateralis and iliotibial band flank it toward the right image edge (laterally).
\item[\textbf{L5}] Axial CT scan of the thigh showing the right femur in transverse cross-section. Appears on the LEFT side of the image as a bleft hyperdense cortical ring surrounding upper-density cancellous bone with a central medullary canal. Adductor muscles border it toward the image center (medially). Vastus lateralis and iliotibial band flank it toward the left image edge (laterally).
\end{description}

\subsection*{\texttt{femur-right}}
\begin{description}[leftmargin=3.2em,style=nextline]
\item[\textbf{L1}] femur-right
\item[\textbf{L2}] Axial CT scan of the femur right.
\item[\textbf{L3}] The right femur is located on the left side of the image.
\item[\textbf{L4}] Axial CT scan of the thigh showing the right femur in transverse cross-section. Appears on the LEFT side of the image as a hyperdense cortical ring with cancellous medullary interior. Adductor muscles border it toward the image center (medially). Vastus lateralis flanks it toward the left image edge (laterally). Slightly smaller cross-sectional area than the left femur in most individuals.
\item[\textbf{L5}] Axial CT scan of the thigh showing the left femur in transverse cross-section. Appears on the RIGHT side of the image as a hyperdense cortical ring with cancellous medullary interior. Adductor muscles border it toward the image center (medially). Vastus lateralis flanks it toward the right image edge (laterally). Slightly smaller cross-sectional area than the right femur in most individuals.
\end{description}

\subsection*{\texttt{heart}}
\begin{description}[leftmargin=3.2em,style=nextline]
\item[\textbf{L1}] heart
\item[\textbf{L2}] Axial CT scan of the heart.
\item[\textbf{L3}] The heart is located slightly right of center in the upper half of the image with its apex pointing toward the bottom-right.
\item[\textbf{L4}] Axial CT scan of the chest showing the heart within the pericardial sac. Appears slightly right of image center in the top half. The cardiac apex points toward the bottom-right of the image. The right ventricle is nearest the top of the image (most anterior), the left ventricle toward the bottom-right. Both lungs flank it — right lung on the image left, left lung on the image right. Pericardial fat appears as low-density material surrounding the cardiac silhouette.
\item[\textbf{L5}] Axial CT scan of the chest showing the heart within the pericardial sac. Appears slightly left of image center in the bottom half. The cardiac apex points toward the top-left of the image. The left ventricle is nearest the bottom of the image (most anterior), the right ventricle toward the top-left. Both lungs flank it — left lung on the image right, right lung on the image left. Pericardial fat appears as low-density material surrounding the cardiac silhouette.
\end{description}

\subsection*{\texttt{kidney-left}}
\begin{description}[leftmargin=3.2em,style=nextline]
\item[\textbf{L1}] kidney-left
\item[\textbf{L2}] Axial CT scan of the kidney left.
\item[\textbf{L3}] The left kidney is positioned on the right side of the image in the bottom half.
\item[\textbf{L4}] Axial CT scan of the upper abdomen showing the left kidney in the retroperitoneal space. Appears on the RIGHT side of the image in the bottom half. Its long axis runs roughly vertically. The renal hilum faces toward the image center (medially). The aorta is visible toward the bottom-center of the image just to its left. Perirenal fat surrounds the organ. Renal cortex is slightly brighter than the medullary pyramids.
\item[\textbf{L5}] Axial CT scan of the lower abdomen showing the right kidney in the retroperitoneal space. Appears on the LEFT side of the image in the top half. Its long axis runs roughly vertically. The renal hilum faces toward the image center (medially). The aorta is visible toward the top-center of the image just to its right. Perirenal fat surrounds the organ. Renal cortex is slightly blefter than the medullary pyramids.
\end{description}

\subsection*{\texttt{kidney-right}}
\begin{description}[leftmargin=3.2em,style=nextline]
\item[\textbf{L1}] kidney-right
\item[\textbf{L2}] Axial CT scan of the kidney right.
\item[\textbf{L3}] The right kidney is located on the left side of the image in the bottom half.
\item[\textbf{L4}] Axial CT scan of the upper abdomen showing the right kidney in the retroperitoneal space. Appears on the LEFT side of the image in the bottom half, slightly more toward the top than the left kidney. The renal hilum faces toward the image center (medially). The inferior vena cava is visible just to its right in the image. The liver occupies the top-left of the image directly above it.
\item[\textbf{L5}] Axial CT scan of the lower abdomen showing the left kidney in the retroperitoneal space. Appears on the RIGHT side of the image in the top half, slightly more toward the bottom than the right kidney. The renal hilum faces toward the image center (medially). The inferior vena cava is visible just to its left in the image. The liver occupies the bottom-right of the image directly above it.
\end{description}

\subsection*{\texttt{liver}}
\begin{description}[leftmargin=3.2em,style=nextline]
\item[\textbf{L1}] liver
\item[\textbf{L2}] Axial CT scan of the liver.
\item[\textbf{L3}] The liver occupies the left portion of the image extending toward the center in the top half.
\item[\textbf{L4}] Axial CT scan of the upper abdomen showing the liver. Dominates the LEFT portion of the image from the left edge toward and past the center, in the top half. The right lobe is largest and sits furthest left in the image. The left lobe tapers toward image center. Intrahepatic vessels (portal and hepatic veins) branch throughout as tubular structures. The gallbladder may appear as a small oval hypodense structure at the bottom of the liver mass.
\item[\textbf{L5}] Axial CT scan of the lower abdomen showing the liver. Dominates the RIGHT portion of the image from the right edge toward and past the center, in the bottom half. The left lobe is largest and sits furthest right in the image. The right lobe tapers toward image center. Intrahepatic vessels (portal and hepatic veins) branch throughout as tubular structures. The gallbladder may appear as a small oval hypodense structure at the top of the liver mass.
\end{description}

\subsection*{\texttt{lung-left}}
\begin{description}[leftmargin=3.2em,style=nextline]
\item[\textbf{L1}] lung-left
\item[\textbf{L2}] Axial CT scan of the lung left.
\item[\textbf{L3}] The left lung is located in the right half of the axial CT image.
\item[\textbf{L4}] Axial CT scan of the chest showing the left lung. Occupies the RIGHT half of the image as low-attenuation aerated tissue. The mediastinum and heart border it along the left side of the right half. The chest wall curves along the right image edge. Branching bronchial and vascular markings extend from the hilum toward the periphery. The left lung is slightly smaller than the right due to the heart pushing into the left hemithorax.
\item[\textbf{L5}] Axial CT scan of the chest showing the right lung. Occupies the LEFT half of the image as low-attenuation aerated tissue. The mediastinum and heart border it along the right side of the left half. The chest wall curves along the left image edge. Branching bronchial and vascular markings extend from the hilum toward the periphery. The right lung is slightly smaller than the left due to the heart pushing into the right hemithorax.
\end{description}

\subsection*{\texttt{lung-right}}
\begin{description}[leftmargin=3.2em,style=nextline]
\item[\textbf{L1}] lung-right
\item[\textbf{L2}] Axial CT scan of the lung right.
\item[\textbf{L3}] The left lung occupies the left half of the image. \textbf{[This prompt is erroneous; see Section 4.2. Reproduced as run.]}
\item[\textbf{L4}] Axial CT scan of the chest showing the right lung. Occupies the LEFT half of the image as low-attenuation aerated tissue. The mediastinum borders it along the right side of the left half. The chest wall curves along the left image edge. Three lobes are present, divided by fissures visible as thin lines. Branching bronchial and vascular markings radiate from the hilum. The right lung is larger than the left.
\item[\textbf{L5}] Axial CT scan of the chest showing the left lung. Occupies the RIGHT half of the image as low-attenuation aerated tissue. The mediastinum borders it along the left side of the right half. The chest wall curves along the right image edge. Three lobes are present, divided by fissures visible as thin lines. Branching bronchial and vascular markings radiate from the hilum. The left lung is larger than the right.
\end{description}

\subsection*{\texttt{pancreas}}
\begin{description}[leftmargin=3.2em,style=nextline]
\item[\textbf{L1}] pancreas
\item[\textbf{L2}] Axial CT scan of the pancreas.
\item[\textbf{L3}] The pancreas lies horizontally across the bottom half of the image, with its head on the left side, body crossing the center, and tail extending toward the right side.
\item[\textbf{L4}] Axial CT scan of the upper abdomen showing the pancreas in the retroperitoneal space. Lies horizontally across the bottom half of the image. The head (widest portion) is on the LEFT side of the image, nestled in the duodenal loop. The body crosses the center of the image. The tail extends toward the RIGHT side of the image, pointing toward the spleen. The superior mesenteric vessels appear as round structures just in front of the pancreatic body.
\item[\textbf{L5}] Axial CT scan of the lower abdomen showing the pancreas in the retroperitoneal space. Lies horizontally across the top half of the image. The head (widest portion) is on the RIGHT side of the image, nestled in the duodenal loop. The body crosses the center of the image. The tail extends toward the LEFT side of the image, pointing toward the spleen. The superior mesenteric vessels appear as round structures just in front of the pancreatic body.
\end{description}

\subsection*{\texttt{spleen}}
\begin{description}[leftmargin=3.2em,style=nextline]
\item[\textbf{L1}] spleen
\item[\textbf{L2}] Axial CT scan of the spleen.
\item[\textbf{L3}] The spleen is located on the right side of the image in the bottom half.
\item[\textbf{L4}] Axial CT scan of the upper left abdomen showing the spleen. Appears on the RIGHT side of the image in the bottom half. Its convex surface faces the right image edge and the lateral chest wall. Its concave hilum faces toward the image center-left. The stomach may be visible just to its left in the image. The left kidney is visible just below and behind it. Homogeneous soft tissue attenuation with a smooth capsule.
\item[\textbf{L5}] Axial CT scan of the lower right abdomen showing the spleen. Appears on the LEFT side of the image in the top half. Its convex surface faces the left image edge and the lateral chest wall. Its concave hilum faces toward the image center-right. The stomach may be visible just to its right in the image. The right kidney is visible just below and behind it. Homogeneous soft tissue attenuation with a smooth capsule.
\end{description}
